% FILE: CSE-25_CT-1.tex
\documentclass{article}
\usepackage{amsmath}
\usepackage{amssymb}
\begin{document}

%%%%%%%%%%%%%%%%%%%%%%%%%%%%%%%%%%%%%%%%%%%%%%%%%%
% QUESTION_START
%%%%%%%%%%%%%%%%%%%%%%%%%%%%%%%%%%%%%%%%%%%%%%%%%%
% id=cse25-ct1-secb-q1i
% discipline=CSE
% year=2025
% exam_type=CT
% section=B
% ct_number=1
% question_number=1i
% topic=Indefinite Integration
% subtopic=Substitution
% difficulty=Hard
% length=Medium
% question_type=New
% frequency=1
% appearances=
% OCR_STATUS=VERIFIED
\section*{CT1 (Sec B) - Q1(i)}
Question:
Evaluate the integral:
\[
\int \sin^{-1} \sqrt{\frac{x}{x+a}} \, dx
\]

Solution:
Step 1: Simplify the integrand using substitution.
Let \(\theta = \sin^{-1} \sqrt{\frac{x}{x+a}}\). This implies:
\[
\sin\theta = \sqrt{\frac{x}{x+a}} \implies \sin^2\theta = \frac{x}{x+a}
\]
\[
(x+a)\sin^2\theta = x \implies a\sin^2\theta = x(1-\sin^2\theta) = x\cos^2\theta
\]
\[
x = a\tan^2\theta
\]

Step 2: Differentiate to find \(dx\).
\[
dx = 2a\tan\theta\sec^2\theta \, d\theta
\]

Step 3: Substitute these into the integral and integrate by parts.
\[
\int \sin^{-1} \sqrt{\frac{x}{x+a}} \, dx = \int \theta \left(2a\tan\theta\sec^2\theta\right) d\theta = 2a \int \theta \left(\tan\theta\sec^2\theta\right) d\theta
\]
Using integration by parts, let:
- \(u = \theta \implies du = d\theta\)
- \(dv = \tan\theta\sec^2\theta \, d\theta \implies v = \frac{1}{2}\tan^2\theta\)

Applying the formula \(\int u \, dv = uv - \int v \, du\):
\[
\int \theta \left(\tan\theta\sec^2\theta\right) d\theta = \frac{1}{2}\theta\tan^2\theta - \frac{1}{2}\int \tan^2\theta \, d\theta
\]
Since \(\int \tan^2\theta \, d\theta = \int (\sec^2\theta - 1) \, d\theta = \tan\theta - \theta\):
\[
\int \theta \left(\tan\theta\sec^2\theta\right) d\theta = \frac{1}{2}\theta\tan^2\theta - \frac{1}{2}(\tan\theta - \theta)
\]

Step 4: Combine the constants and substitute back to \(x\).
Multiplying by \(2a\):
\[
2a \left[ \frac{1}{2}\theta\tan^2\theta - \frac{1}{2}(\tan\theta - \theta) \right] = a\theta\tan^2\theta - a\tan\theta + a\theta = a\theta(1+\tan^2\theta) - a\tan\theta
\]
Recall that \(\tan^2\theta = \frac{x}{a}\), \(\tan\theta = \sqrt{\frac{x}{a}}\), and \(\theta = \tan^{-1}\sqrt{\frac{x}{a}}\):
\[
= a\left(\tan^{-1}\sqrt{\frac{x}{a}}\right)\left(1+\frac{x}{a}\right) - a\sqrt{\frac{x}{a}}
\]
\[
= (x+a)\tan^{-1}\sqrt{\frac{x}{a}} - \sqrt{ax} + C
\]

Final Answer:
\[
\boxed{(x+a)\tan^{-1}\sqrt{\frac{x}{a}} - \sqrt{ax} + C}
\]
%%%%%%%%%%%%%%%%%%%%%%%%%%%%%%%%%%%%%%%%%%%%%%%%%%
% QUESTION_END
%%%%%%%%%%%%%%%%%%%%%%%%%%%%%%%%%%%%%%%%%%%%%%%%%%

%%%%%%%%%%%%%%%%%%%%%%%%%%%%%%%%%%%%%%%%%%%%%%%%%%
% QUESTION_START
%%%%%%%%%%%%%%%%%%%%%%%%%%%%%%%%%%%%%%%%%%%%%%%%%%
% id=cse25-ct1-secb-q1ii
% discipline=CSE
% year=2025
% exam_type=CT
% section=B
% ct_number=1
% question_number=1ii
% topic=Indefinite Integration
% subtopic=Substitution
% difficulty=Medium
% length=Medium
% question_type=Repeated PYQ Pattern
% frequency=1
% appearances=
% OCR_STATUS=VERIFIED
\section*{CT1 (Sec B) - Q1(ii)}
Question:
Evaluate the integral:
\[
\int \frac{4e^x+6e^{-x}}{9e^x-4e^{-x}} \, dx
\]

Solution:
Step 1: Express the numerator in terms of the denominator and its derivative.
Let the numerator equal \(A \cdot (\text{Denominator}) + B \cdot (\text{Derivative of Denominator})\):
\[
4e^x + 6e^{-x} = A(9e^x - 4e^{-x}) + B \frac{d}{dx}(9e^x - 4e^{-x})
\]
Since \(\frac{d}{dx}(9e^x - 4e^{-x}) = 9e^x + 4e^{-x}\):
\[
4e^x + 6e^{-x} = A(9e^x - 4e^{-x}) + B(9e^x + 4e^{-x})
\]

Step 2: Equate coefficients to solve for \(A\) and \(B\).
For the coefficients of \(e^x\):
\[
4 = 9A + 9B \implies A + B = \frac{4}{9} \quad \text{--- (Equation 1)}
\]
For the coefficients of \(e^{-x}\):
\[
6 = -4A + 4B \implies -A + B = \frac{3}{2} \quad \text{--- (Equation 2)}
\]
Adding Equations 1 and 2:
\[
2B = \frac{4}{9} + \frac{3}{2} = \frac{8+27}{18} = \frac{35}{18} \implies B = \frac{35}{36}
\]
Subtracting Equation 2 from Equation 1:
\[
2A = \frac{4}{9} - \frac{3}{2} = \frac{8-27}{18} = -\frac{19}{18} \implies A = -\frac{19}{36}
\]

Step 3: Integrate using the decomposed terms.
We can rewrite the integral as:
\[
\int \frac{4e^x + 6e^{-x}}{9e^x - 4e^{-x}} \, dx = \int \left( A + B \frac{9e^x + 4e^{-x}}{9e^x - 4e^{-x}} \right) dx
\]
\[
= Ax + B \ln|9e^x - 4e^{-x}| + C
\]
Substituting the values of \(A\) and \(B\):
\[
= -\frac{19}{36}x + \frac{35}{36}\ln|9e^x - 4e^{-x}| + C
\]

Final Answer:
\[
\boxed{-\frac{19}{36}x + \frac{35}{36}\ln|9e^x - 4e^{-x}| + C}
\]
%%%%%%%%%%%%%%%%%%%%%%%%%%%%%%%%%%%%%%%%%%%%%%%%%%
% QUESTION_END
%%%%%%%%%%%%%%%%%%%%%%%%%%%%%%%%%%%%%%%%%%%%%%%%%%

%%%%%%%%%%%%%%%%%%%%%%%%%%%%%%%%%%%%%%%%%%%%%%%%%%
% QUESTION_START
%%%%%%%%%%%%%%%%%%%%%%%%%%%%%%%%%%%%%%%%%%%%%%%%%%
% id=cse25-ct1-secb-q2
% discipline=CSE
% year=2025
% exam_type=CT
% section=B
% ct_number=1
% question_number=2
% topic=Definite Integration
% subtopic=Reduction Formula
% difficulty=Medium
% length=Medium
% question_type=Repeated PYQ Pattern
% frequency=1
% appearances=
% OCR_STATUS=VERIFIED
\section*{CT1 (Sec B) - Q2}
Question:
Obtain a reduction formula for \(\int \sin^m x \, dx\) and apply it to evaluate \(\int_0^{\pi/2} \sin^9 x \, dx\).

Solution:
Step 1: Set up the reduction formula for the indefinite integral.
Let \(I_m = \int \sin^m x \, dx = \int \sin^{m-1} x \sin x \, dx\).
Applying integration by parts:
- Let \(u = \sin^{m-1} x \implies du = (m-1)\sin^{m-2} x \cos x \, dx\)
- Let \(dv = \sin x \, dx \implies v = -\cos x\)

\[
I_m = -\sin^{m-1} x \cos x + (m-1) \int \sin^{m-2} x \cos^2 x \, dx
\]
Substitute \(\cos^2 x = 1 - \sin^2 x\):
\[
I_m = -\sin^{m-1} x \cos x + (m-1) \int \sin^{m-2} x (1 - \sin^2 x) \, dx
\]
\[
I_m = -\sin^{m-1} x \cos x + (m-1) I_{m-2} - (m-1) I_m
\]
\[
m I_m = -\sin^{m-1} x \cos x + (m-1) I_{m-2}
\]
Thus, the reduction formula is:
\[
I_m = -\frac{\sin^{m-1} x \cos x}{m} + \frac{m-1}{m} I_{m-2}
\]

Step 2: Apply the reduction formula to the definite integral.
Let \(J_m = \int_0^{\pi/2} \sin^m x \, dx\). The boundary term evaluated at \(0\) and \(\pi/2\) is \(0\). Thus, we get:
\[
J_m = \frac{m-1}{m} J_{m-2}
\]

Step 3: Evaluate \(J_9\).
\[
J_9 = \frac{8}{9} J_7 = \frac{8}{9} \cdot \frac{6}{7} J_5 = \frac{8}{9} \cdot \frac{6}{7} \cdot \frac{4}{5} J_3 = \frac{8}{9} \cdot \frac{6}{7} \cdot \frac{4}{5} \cdot \frac{2}{3} J_1
\]
Since \(J_1 = \int_0^{\pi/2} \sin x \, dx = [-\cos x]_0^{\pi/2} = 1\):
\[
J_9 = \frac{8 \cdot 6 \cdot 4 \cdot 2}{9 \cdot 7 \cdot 5 \cdot 3} = \frac{384}{945} = \frac{128}{315}
\]

Final Answer:
\[
\boxed{\frac{128}{315}}
\]
%%%%%%%%%%%%%%%%%%%%%%%%%%%%%%%%%%%%%%%%%%%%%%%%%%
% QUESTION_END
%%%%%%%%%%%%%%%%%%%%%%%%%%%%%%%%%%%%%%%%%%%%%%%%%%

%%%%%%%%%%%%%%%%%%%%%%%%%%%%%%%%%%%%%%%%%%%%%%%%%%
% QUESTION_START
%%%%%%%%%%%%%%%%%%%%%%%%%%%%%%%%%%%%%%%%%%%%%%%%%%
% id=cse25-ct1-secb-q3
% discipline=CSE
% year=2025
% exam_type=CT
% section=B
% ct_number=1
% question_number=3
% topic=Definite Integration
% subtopic=General
% difficulty=Medium
% length=Medium
% question_type=New
% frequency=1
% appearances=
% OCR_STATUS=VERIFIED
\section*{CT1 (Sec B) - Q3}
Question:
Applying the definition of definite integral, prove that the limit of the sum of the series \(\frac{1}{n} + \frac{1}{n+1} + \frac{1}{n+2} + \dots + \frac{1}{3n} = \log 3\).

Solution:
Step 1: Write the sum in summation notation.
Let \(S_n\) be the given series. Note that \(\frac{1}{3n} = \frac{1}{n+2n}\).
\[
S_n = \sum_{r=0}^{2n} \frac{1}{n+r}
\]

Step 2: Convert the sum to a Riemann sum format.
Factor out \(1/n\) from each term:
\[
S_n = \frac{1}{n} \sum_{r=0}^{2n} \frac{1}{1 + \frac{r}{n}}
\]

Step 3: Convert the limit of sum to a definite integral.
By definition, as \(n \to \infty\), let \(\frac{r}{n} \to x\) and \(\frac{1}{n} \to dx\).
The limits of integration are:
- For \(r = 0 \implies x = 0\)
- For \(r = 2n \implies x = 2\)
Thus:
\[
\lim_{n \to \infty} S_n = \int_0^2 \frac{1}{1+x} \, dx
\]

Step 4: Integrate to prove the limit value.
\[
\int_0^2 \frac{1}{1+x} \, dx = \left[ \log(1+x) \right]_0^2 = \log 3 - \log 1 = \log 3
\]
Hence proved.

Final Answer:
\[
\boxed{\log 3}
\]
%%%%%%%%%%%%%%%%%%%%%%%%%%%%%%%%%%%%%%%%%%%%%%%%%%
% QUESTION_END
%%%%%%%%%%%%%%%%%%%%%%%%%%%%%%%%%%%%%%%%%%%%%%%%%%

%%%%%%%%%%%%%%%%%%%%%%%%%%%%%%%%%%%%%%%%%%%%%%%%%%
% QUESTION_START
%%%%%%%%%%%%%%%%%%%%%%%%%%%%%%%%%%%%%%%%%%%%%%%%%%
% id=cse25-ct1-secb-q4
% discipline=CSE
% year=2025
% exam_type=CT
% section=B
% ct_number=1
% question_number=4
% topic=Definite Integration
% subtopic=General
% difficulty=Medium
% length=Long
% question_type=New
% frequency=1
% appearances=
% OCR_STATUS=VERIFIED
\section*{CT1 (Sec B) - Q4}
Question:
State the key properties of definite integrals. Using these properties, evaluate \(\int_0^{\pi/2} \frac{\sqrt{\tan x}}{1 + \sqrt{\tan x}} \, dx\).

Solution:
Step 1: State key properties of definite integrals.
1. \(\int_a^b f(x) \, dx = \int_a^b f(t) \, dt\)
2. \(\int_a^b f(x) \, dx = -\int_b^a f(x) \, dx\)
3. \(\int_a^b [f(x) \pm g(x)] \, dx = \int_a^b f(x) \, dx \pm \int_a^b g(x) \, dx\)
4. \(\int_a^b f(x) \, dx = \int_a^c f(x) \, dx + \int_c^b f(x) \, dx\)
5. \(\int_a^b f(x) \, dx = \int_a^b f(a+b-x) \, dx\)
6. \(\int_{-a}^a f(x) \, dx = 2\int_0^a f(x) \, dx\) (if \(f\) is even), or \(0\) (if \(f\) is odd).

Step 2: Simplify the integrand using sine and cosine.
Let \(I = \int_0^{\pi/2} \frac{\sqrt{\tan x}}{1 + \sqrt{\tan x}} \, dx\).
Substitute \(\tan x = \frac{\sin x}{\cos x}\):
\[
I = \int_0^{\pi/2} \frac{\sqrt{\sin x}}{\sqrt{\cos x} + \sqrt{\sin x}} \, dx \quad \text{--- (Equation 1)}
\]

Step 3: Apply the property \(\int_0^a f(x) \, dx = \int_0^a f(a-x) \, dx\).
\[
I = \int_0^{\pi/2} \frac{\sqrt{\sin\left(\frac{\pi}{2}-x\right)}}{\sqrt{\cos\left(\frac{\pi}{2}-x\right)} + \sqrt{\sin\left(\frac{\pi}{2}-x\right)}} \, dx
\]
\[
I = \int_0^{\pi/2} \frac{\sqrt{\cos x}}{\sqrt{\sin x} + \sqrt{\cos x}} \, dx \quad \text{--- (Equation 2)}
\]

Step 4: Sum Equation 1 and Equation 2.
\[
2I = \int_0^{\pi/2} \frac{\sqrt{\sin x} + \sqrt{\cos x}}{\sqrt{\sin x} + \sqrt{\cos x}} \, dx
\]
\[
2I = \int_0^{\pi/2} 1 \, dx = \left[ x \right]_0^{\pi/2} = \frac{\pi}{2}
\]
\[
I = \frac{\pi}{4}
\]

Final Answer:
\[
\boxed{\frac{\pi}{4}}
\]
%%%%%%%%%%%%%%%%%%%%%%%%%%%%%%%%%%%%%%%%%%%%%%%%%%
% QUESTION_END
%%%%%%%%%%%%%%%%%%%%%%%%%%%%%%%%%%%%%%%%%%%%%%%%%%

%%%%%%%%%%%%%%%%%%%%%%%%%%%%%%%%%%%%%%%%%%%%%%%%%%
% QUESTION_START
%%%%%%%%%%%%%%%%%%%%%%%%%%%%%%%%%%%%%%%%%%%%%%%%%%
% id=cse25-ct1-secb-q5i
% discipline=CSE
% year=2025
% exam_type=CT
% section=B
% ct_number=1
% question_number=5i
% topic=Definite Integration
% subtopic=General
% difficulty=Hard
% length=Long
% question_type=Repeated PYQ Pattern
% frequency=1
% appearances=
% OCR_STATUS=VERIFIED
\section*{CT1 (Sec B) - Q5(i)}
Question:
Using the properties of the definite integral, integrate:
\[
\int_0^{\pi/2} \log(\sin x) \, dx
\]

Solution:
Step 1: Set up the integral with the definite property.
Let \(I = \int_0^{\pi/2} \log(\sin x) \, dx\).
Using \(\int_0^a f(x) \, dx = \int_0^a f(a-x) \, dx\):
\[
I = \int_0^{\pi/2} \log\left(\sin\left(\frac{\pi}{2}-x\right)\right) \, dx = \int_0^{\pi/2} \log(\cos x) \, dx
\]

Step 2: Add both expressions for \(I\).
\[
2I = \int_0^{\pi/2} \log(\sin x) \, dx + \int_0^{\pi/2} \log(\cos x) \, dx = \int_0^{\pi/2} \log(\sin x \cos x) \, dx
\]
Substitute \(\sin x \cos x = \frac{\sin 2x}{2}\):
\[
2I = \int_0^{\pi/2} \log\left(\frac{\sin 2x}{2}\right) \, dx = \int_0^{\pi/2} \log(\sin 2x) \, dx - \int_0^{\pi/2} \log 2 \, dx
\]
\[
2I = \int_0^{\pi/2} \log(\sin 2x) \, dx - \frac{\pi}{2}\log 2 \quad \text{--- (Equation 1)}
\]

Step 3: Use substitution to simplify the remaining integral.
Let \(J = \int_0^{\pi/2} \log(\sin 2x) \, dx\).
Substitute \(t = 2x \implies dx = \frac{dt}{2}\).
- When \(x = 0 \implies t = 0\)
- When \(x = \pi/2 \implies t = \pi\)
\[
J = \frac{1}{2} \int_0^{\pi} \log(\sin t) \, dt
\]
Using the symmetry property \(\int_0^{2a} f(t) \, dt = 2\int_0^a f(t) \, dt\) because \(\sin(\pi-t) = \sin t\):
\[
J = \frac{1}{2} \left( 2 \int_0^{\pi/2} \log(\sin t) \, dt \right) = \int_0^{\pi/2} \log(\sin t) \, dt = I
\]

Step 4: Substitute back into Equation 1 and solve for \(I\).
\[
2I = I - \frac{\pi}{2}\log 2 \implies I = -\frac{\pi}{2}\log 2 = \frac{\pi}{2}\log\left(\frac{1}{2}\right)
\]

Final Answer:
\[
\boxed{\frac{\pi}{2}\log\left(\frac{1}{2}\right)}
\]
%%%%%%%%%%%%%%%%%%%%%%%%%%%%%%%%%%%%%%%%%%%%%%%%%%
% QUESTION_END
%%%%%%%%%%%%%%%%%%%%%%%%%%%%%%%%%%%%%%%%%%%%%%%%%%

%%%%%%%%%%%%%%%%%%%%%%%%%%%%%%%%%%%%%%%%%%%%%%%%%%
% QUESTION_START
%%%%%%%%%%%%%%%%%%%%%%%%%%%%%%%%%%%%%%%%%%%%%%%%%%
% id=cse25-ct1-secb-q5ii
% discipline=CSE
% year=2025
% exam_type=CT
% section=B
% ct_number=1
% question_number=5ii
% topic=Definite Integration
% subtopic=General
% difficulty=Medium
% length=Medium
% question_type=New
% frequency=1
% appearances=
% OCR_STATUS=VERIFIED
\section*{CT1 (Sec B) - Q5(ii)}
Question:
Using the properties of the definite integral, integrate:
\[
\int_{-3}^3 |x - 1| \, dx
\]

Solution:
Step 1: Identify the transition point of the absolute value.
The absolute value term changes behavior at \(x - 1 = 0 \implies x = 1\).
- For \(x \in [-3, 1]\), \(|x - 1| = -(x - 1) = 1 - x\).
- For \(x \in [-3, 1]\), \(|x - 1| = -(x - 1) = 1 - x\).
Wait, there is a typo in my thoughts. Over the interval $[-3, 1]$, $|x-1| = -(x-1) = 1-x$.
For $x \in [1, 3]$, $|x-1| = x-1$.
The split is correct.

Step 2: Split the integral over the subintervals.
\[
\int_{-3}^3 |x - 1| \, dx = \int_{-3}^1 (1 - x) \, dx + \int_1^3 (x - 1) \, dx
\]

Step 3: Integrate and evaluate both parts.
- First part:
  \[
  \int_{-3}^1 (1 - x) \, dx = \left[ x - \frac{x^2}{2} \right]_{-3}^1 = \left(1 - \frac{1}{2}\right) - \left(-3 - \frac{9}{2}\right) = \frac{1}{2} - \left(-\frac{15}{2}\right) = 8
  \]
- Second part:
  \[
  \int_1^3 (x - 1) \, dx = \left[ \frac{x^2}{2} - x \right]_1^3 = \left(\frac{9}{2} - 3\right) - \left(\frac{1}{2} - 1\right) = \frac{3}{2} - \left(-\frac{1}{2}\right) = 2
  \]

Step 4: Sum both evaluated values.
\[
\int_{-3}^3 |x - 1| \, dx = 8 + 2 = 10
\]

Final Answer:
\[
\boxed{10}
\]
%%%%%%%%%%%%%%%%%%%%%%%%%%%%%%%%%%%%%%%%%%%%%%%%%%
% QUESTION_END
%%%%%%%%%%%%%%%%%%%%%%%%%%%%%%%%%%%%%%%%%%%%%%%%%%

\end{document}